\documentclass{article}

\usepackage[english]{babel}
\usepackage[letterpaper,top=2cm,bottom=2cm,left=2.5cm,right=2.5cm,marginparwidth=1.75cm]{geometry}

\usepackage{amsmath}
\usepackage{graphicx}
\usepackage{booktabs}
\usepackage{subcaption}
\usepackage{xspace}
\usepackage[colorlinks=true, allcolors=blue]{hyperref}

% Plot search paths (one per evidence subfolder)
\graphicspath{%
  {evidence/01_seed_variance/plots/}%
  {evidence/02_canon_performance/plots/}%
  {evidence/03_canon_ablations/plots/}%
  {evidence/04_training_dynamics/plots/}%
  {evidence/05_task_mixing_variance/plots/}%
  {evidence/06_scaling_laws/plots/}%
  {evidence/07_dynamics_on_synthetic/plots/}%
  {evidence/08_canon_stability/plots/}%
  {evidence/09_seed_flip/plots/}%
  {evidence/10_variance_recompute/plots/}%
}

% --- macros ---
\newcommand{\Canon}{\textsc{Canon}\xspace}
\newcommand{\Depo}{\textsc{Depo}\xspace}
\newcommand{\Brevo}{\textsc{Brevo}\xspace}
\newcommand{\Bfs}{\textsc{Bfs}\xspace}
\newcommand{\Concomp}{\textsc{ConComp}\xspace}
\newcommand{\ShortPath}{\textsc{ShortPath}\xspace}
\newcommand{\Llama}{\textsc{Llama}\xspace}

\title{Trustworthy Verdicts on Synthetic Reasoning Benchmarks:\\
A Reliability Protocol and a Conditional Re-Evaluation of the Canon Layer}
\author{Kirill Zemlianskii}
\date{\today}

\begin{document}
\maketitle

\begin{abstract}
Synthetic reasoning benchmarks promise a cheap, controlled testbed for validating
architectural changes to transformers. We show that this promise is threatened on two
compounding axes: \emph{seed variance} corrupts the measurement, and \emph{aggregation}
hides where an effect actually lives. Using the recently proposed \Canon layer (a learned
causal 1-D convolution inserted within each transformer block) as a running case study, we
first characterize the seed-variance failure --- the same architecture grokks on some seeds
and stagnates on others, while text modeling is stable --- and examine two levers commonly
assumed to fix it. Task mixing reduces the \emph{aggregate} seed variance, but partly by
driving the hard task to a reliably-failed floor; \Canon's lower per-task spread is mostly a
mean effect (the coefficient of variation is essentially unchanged). The climax of the
reliability analysis is a seed-flip audit: pooling all sixteen available 30M / 8-task seeds,
the headline \Canon-vs-\Llama verdict is \emph{statistically null} (gap $+0.015$, $\sim$1.3$\sigma$,
permutation $p=0.22$) and its \emph{sign flips} across seed batches. A single aggregate number
at a single scale is therefore the wrong object of study. We then re-evaluate \Canon
conditionally and find its benefit is real but scoped: large and seed-robust at small scale
and on multi-hop retrieval (\Depo), null in the 30M aggregate, and associated with control of
deep-layer gradient dominance. Our primary contribution is a multi-seed, multi-task,
multi-scale evaluation protocol for benchmarks of this kind, demonstrated on a verdict that
flips under exactly the single-number protocol the field currently uses.
\end{abstract}

% ===========================================================================
\section{Introduction}
\label{sec:intro}
% ===========================================================================

Experimenting with architectural changes on real language data is expensive, so the field
increasingly turns to synthetic reasoning benchmarks --- controlled graph and sequence tasks
that isolate a specific transformer capability and train to convergence in hours rather than
weeks. The Physics of Language Models framework~\cite{allen_zhu_physics4} exemplifies this
program: it provides both a battery of such tasks and a candidate architectural improvement,
the \Canon layer, a learned causal 1-D convolution inserted at one or more positions inside
each transformer block. The original work reports that \Canon improves performance and leaves
open why it helps, which hyperparameters matter, and whether the improvement is consistent
across scale.

A benchmark, however, is only as useful as the trustworthiness of the verdicts it produces.
We find that verdicts on these benchmarks are threatened on two distinct axes. The first is a
\emph{measurement} problem: the tasks exhibit extreme seed-to-seed variance --- the same
architecture grokks on some random seeds and never escapes the initial plateau on others ---
so a comparison between two architectures run at one seed each can report either sign. The
second is a \emph{validity} problem: even with the noise tamed, a single aggregate score across
a task mixture conflates tasks at the accuracy ceiling, tasks at the floor, and the one or two
tasks where an architectural effect actually concentrates. Both axes share one villain: being
fooled by a benchmark number.

We make these threats concrete by threading \Canon through the entire paper as a running
example. The irony that powers our study is that the headline comparison nearly everyone would
report --- ``\Canon beats \Llama at 30M parameters on the standard 8-task mix'' --- is, measured
under a protocol that uses every seed we have, \emph{null}.

\paragraph{Contributions.}
\begin{itemize}
  \item We characterize the seed-variance failure mode of Physics-style synthetic benchmarks
        and quantify two commonly-assumed variance fixes (task mixing and architecture),
        showing that each is partly an artifact: a floor effect and a mean effect respectively
        (Section~\ref{sec:p1}).
  \item We introduce a seed-flip audit and show that the 30M / 8-task \Canon-vs-\Llama verdict
        is statistically null (permutation $p=0.22$) and flips by seed batch, motivating a
        multi-seed, multi-task, multi-scale evaluation protocol that reports verdict stability
        as a function of seed count (Section~\ref{sec:p1:flip}).
  \item Under that protocol we re-evaluate \Canon conditionally: its advantage is robust only
        at small scale and on multi-hop retrieval (\Depo), dissolves into seed noise by 30M,
        and is mechanistically associated with suppression of deep-layer gradient dominance
        (Section~\ref{sec:p2}).
\end{itemize}

To make the threats precise we first define the tasks and the \Canon layer.

% ===========================================================================
\section{Background and Problem Statement}
\label{sec:background}
% ===========================================================================

\subsection{Synthetic Reasoning Tasks}
\label{sec:background:tasks}

Each task encodes a graph or sequence problem as a token stream and trains the model to
predict only the answer tokens. Two encodings are used throughout: \texttt{edges\_list}
(the graph as an unordered edge set) and \texttt{adj\_list} (grouped by source node).
Difficulty is governed by knobs such as the maximum number of nodes and the maximum number
of hops.

\paragraph{\Depo} (multi-hop retrieval) encodes the $n$ edges of a circular permutation
followed by a query $(q,k)$; the model must emit the $k$-th successor $f^{k}(q)$ with no
intermediate chain of thought.
\paragraph{\Concomp} (connected-component factorization) requires partitioning a graph into
its connected components; we score prefix accuracy.
\paragraph{\ShortPath} requires emitting a shortest path between two vertices; we score whether
the returned path is correct.
\paragraph{\Bfs} requires emitting the set of vertices reachable from a source; we score set
recall. \Brevo (reasoning breadth over a DAG) is part of the original battery but is excluded
from the standard mixture due to extreme instability.

The \emph{standard mixture} used in our scaling experiments is the \textbf{8-task mix}: the four
tasks above under both encodings. Performance on the mixture is summarized by the average
next-token loss on held-out answer tokens (lower is better).

\subsection{The Canon Layer}
\label{sec:background:canon}

A \Canon layer is a learned linear convolution over the last $k$ hidden states along the
sequence dimension. With hidden state $\mathbf{h}_t$ at position $t$ and kernel size $k$,
the output at position $t$ is
\begin{equation}
  \tilde{\mathbf{h}}_t = \sum_{i=0}^{k-1} w_i\, \mathbf{h}_{t-i},
  \label{eq:canon}
\end{equation}
where the $\{w_i\}$ are scalar weights shared across the feature dimension; the convolution is
causal, so position $t$ never sees the future. A \Canon layer can be inserted at four positions
within each transformer block: \textbf{A} (pre-attention), \textbf{B} (post-attention),
\textbf{C} (pre-FFN), and \textbf{D} (post-FFN). Unless stated otherwise we use \Canon-ABCD
(all four positions, all layers) with kernel size $k=4$. The hyperparameters of interest are
$k$, the initialization of $\{w_i\}$ in~\eqref{eq:canon}, and the subset of positions and
layers used.

\subsection{Two Threats to Validity}
\label{sec:background:threats}

We state the two problems that the rest of the paper addresses.
\textbf{(T1) Measurement.} Synthetic-task training is seed-treacherous: the variance of the
final score across random seeds can be larger than the architectural effect being measured.
\textbf{(T2) Validity.} A single aggregate score over a task mixture is a weighted average that
is dominated by easy (near-ceiling) tasks and can mask both floor effects and the localization
of an effect to one or two tasks.

We address T1 first: can we even measure a difference reliably?

% ===========================================================================
\section{Part 1: Can We Trust the Measurement?}
\label{sec:p1}
% ===========================================================================

\subsection{Seed Variance Is Real and Task-Specific}
\label{sec:p1:variance}

\begin{figure*}[t]
  \centering
  \begin{subfigure}[b]{0.32\linewidth}
    \includegraphics[width=\linewidth]{text_stable.png}
    \caption{Text modeling (stable).}
    \label{fig:f1:text}
  \end{subfigure}\hfill
  \begin{subfigure}[b]{0.32\linewidth}
    \includegraphics[width=\linewidth]{depo_8l_noisy.png}
    \caption{\Depo, 8 layers (noisy).}
    \label{fig:f1:depo8}
  \end{subfigure}\hfill
  \begin{subfigure}[b]{0.32\linewidth}
    \includegraphics[width=\linewidth]{depo_12l_grokking.png}
    \caption{\Depo, 12 layers (grokking).}
    \label{fig:f1:depo12}
  \end{subfigure}
  \caption{\textbf{The benchmark used to validate architectures is seed-treacherous.}
    On text modeling (a) the seed-to-seed spread in final loss is below $0.002$. On \Depo the
    same architecture either converges far faster on one seed than the others (b) or groks
    after a long plateau on one seed while another never solves the task (c). Variance this
    large can by itself flip an architectural verdict.}
  \label{fig:f1}
\end{figure*}

Training the same architecture under different random seeds on \Depo and \Brevo produces
qualitatively different loss trajectories --- grokking on some seeds, permanent stagnation on
others --- whereas standard text modeling on FineWeb-Edu yields nearly identical curves
(Figure~\ref{fig:f1}). The seed-to-seed final-loss spread is below $0.002$ on text but reaches
$\approx 0.11$ on \Depo. The instability is therefore task-specific, not generic optimizer
noise: it is a property of the discrete reasoning objective, and it is large enough to flip a
single-seed comparison between two architectures. If we want to compare \Canon against \Llama,
we must reduce this variance before trusting any number. Two levers are available.

\subsection{Two Variance Levers --- and What They Hide}
\label{sec:p1:levers}

\begin{figure*}[t]
  \centering
  \begin{subfigure}[b]{0.32\linewidth}
    \includegraphics[width=\linewidth]{std_vs_tasks.png}
    \caption{Aggregate std vs.\ \#tasks.}
    \label{fig:f2:agg}
  \end{subfigure}\hfill
  \begin{subfigure}[b]{0.32\linewidth}
    \includegraphics[width=\linewidth]{f5_per_task_variance.png}
    \caption{Per-task \Depo std (floor effect).}
    \label{fig:f2:floor}
  \end{subfigure}\hfill
  \begin{subfigure}[b]{0.32\linewidth}
    \includegraphics[width=\linewidth]{f7_cv.png}
    \caption{\Canon vs.\ \Llama on \Depo$\times2$: CV.}
    \label{fig:f2:cv}
  \end{subfigure}
  \caption{\textbf{Both variance levers are partly artifacts.}
    (a) Mixing more tasks lowers the aggregate seed std. (b) But the fixed-task (\Depo) seed std
    only collapses at 8 tasks because \Depo accuracy craters to $0.045$ --- the aggregate
    stabilizes partly by driving the hard task to a reliably-failed floor, not by solving it.
    (c) \Canon's raw std on \Depo$\times2$ is 46\% below \Llama's, but its coefficient of
    variation is essentially identical ($0.166$ vs.\ $0.181$): the spread reduction follows the
    lower mean, not an independent variance-reduction mechanism.}
  \label{fig:f2}
\end{figure*}

\paragraph{Data-side: task mixing.} Training on more diverse tasks reduces the aggregate seed
std monotonically, from $0.045$ on a single task to $0.011$ on the 8-task mix
(Figure~\ref{fig:f2:agg}). This is real, but the per-task picture is less flattering: the
seed std of the fixed hard task (\Depo) only collapses at the 8-task mix because \Depo
accuracy itself craters to $0.045$ (Figure~\ref{fig:f2:floor}). Mixing stabilizes the
\emph{aggregate} partly by pushing the hard task to a reliably-failed floor rather than by
genuinely stabilizing its training. The aggregate is also a mean across heterogeneous tasks,
so its absolute level should not be compared across mixtures of different composition.

\paragraph{Architecture-side: \Canon.} On \Depo$\times2$ (five seeds each), \Canon's raw final-loss
std is $0.039$ versus \Llama's $0.072$ --- a 46\% reduction that invites the claim that \Canon
is an architectural variance-reduction mechanism. But \Canon also reaches a much lower mean loss
($0.237$ vs.\ $0.401$), and once we normalize, the coefficient of variation is essentially equal
($0.166$ vs.\ $0.181$, only 8\% lower; Figure~\ref{fig:f2:cv}). The spread reduction tracks the
lower mean; it is not an independent stabilizing mechanism.

Neither lever makes a single aggregate number trustworthy. The question that decides the
reliability of the whole comparison is whether the headline verdict survives once we use every
seed available.

\subsection{The Residual Problem: the 30M Verdict Is Null and Flips}
\label{sec:p1:flip}

\begin{figure*}[t]
  \centering
  \begin{subfigure}[b]{0.32\linewidth}
    \includegraphics[width=\linewidth]{per_seed_loss.png}
    \caption{Per-seed final loss.}
    \label{fig:f3:seed}
  \end{subfigure}\hfill
  \begin{subfigure}[b]{0.32\linewidth}
    \includegraphics[width=\linewidth]{batch_flip.png}
    \caption{Verdict flips by seed batch.}
    \label{fig:f3:batch}
  \end{subfigure}\hfill
  \begin{subfigure}[b]{0.32\linewidth}
    \includegraphics[width=\linewidth]{verdict_vs_seedcount.png}
    \caption{Verdict stability vs.\ seed count.}
    \label{fig:f3:count}
  \end{subfigure}
  \caption{\textbf{The headline 30M / 8-task verdict is null and seed-fragile.}
    Pooling all 8 \Canon and 8 \Llama seeds, the mean gap is $+0.015$ ($\sim$1.3$\sigma$),
    permutation $p=0.22$ (a). The sign of the verdict flips by seed batch (b): the dynamics
    batch favors \Llama ($0.176$ vs.\ $0.180$), the scaling batch favors \Canon ($0.173$ vs.\
    $0.198$). Even the sign is unstable: at one seed the ``correct'' sign appears 67\% of the
    time, and at eight seeds only 91\%, with the 95\% interval on the gap still including zero
    (c).}
  \label{fig:f3}
\end{figure*}

We pool every available 30M / 8-task run: eight \Canon seeds and eight \Llama seeds, drawn from
two independently-launched batches (a dynamics batch and a scaling batch). The aggregate mean
gap is $+0.015$ in \Canon's favor --- about $1.3\sigma$ --- and a permutation test on the seed
labels gives $p=0.22$: \textbf{not significant} (Figure~\ref{fig:f3:seed}). Worse, the verdict's
\emph{sign} is not stable. Split by batch, the dynamics batch reports \Llama better ($0.176$ vs.\
$0.180$) while the scaling batch reports \Canon better ($0.173$ vs.\ $0.198$); the two batches
disagree on who wins (Figure~\ref{fig:f3:batch}). Resampling the verdict as a function of seed
count (Figure~\ref{fig:f3:count}) shows that a single seed recovers the pooled-mean sign only
67\% of the time, and even eight seeds reach only 91\%, with a 95\% confidence interval on the
gap that still includes zero.

The headline that nearly everyone would report --- ``\Canon wins at 30M on the standard mix''
--- is, measured honestly, null. A single aggregate number at a single scale is the wrong object
of study: it hides seed noise on top of the ceiling and floor effects of
Section~\ref{sec:p1:levers}. We therefore reframe the question from ``does \Canon win'' to
``at which scales and on which tasks is there a robust effect, and why'' --- questions that
remain answerable precisely where the aggregate is not.

% ===========================================================================
\section{Part 2: When and Why Does Canon Help?}
\label{sec:p2}
% ===========================================================================

Since the 30M aggregate carries no signal (Section~\ref{sec:p1:flip}), the positive case for
\Canon must be made where a robust effect actually exists. We proceed in three steps: which
configuration matters (Section~\ref{sec:p2:config}), where the gains live across scale and task
(Section~\ref{sec:p2:where}), and what mechanism underlies them
(Section~\ref{sec:p2:mechanism}).

\subsection{What Configuration Matters}
\label{sec:p2:config}

\begin{figure*}[t]
  \centering
  \begin{subfigure}[b]{0.32\linewidth}
    \includegraphics[width=\linewidth]{kernel_size.png}
    \caption{Kernel size.}
    \label{fig:f4:k}
  \end{subfigure}\hfill
  \begin{subfigure}[b]{0.32\linewidth}
    \includegraphics[width=\linewidth]{initialization.png}
    \caption{Initialization.}
    \label{fig:f4:init}
  \end{subfigure}\hfill
  \begin{subfigure}[b]{0.32\linewidth}
    \includegraphics[width=\linewidth]{layer_depth.png}
    \caption{Layer depth.}
    \label{fig:f4:depth}
  \end{subfigure}
  \caption{\textbf{Configuration ablations on the stable text signal.} Kernel size: $k=1$
    (pointwise rescaling) is worse than no \Canon, and performance improves monotonically for
    $k\ge2$, so cross-token mixing is the essential ingredient. Initialization: every scheme,
    including zero init, beats the baseline. Layer depth: applying \Canon to all layers is best;
    first-layer-only still beats the baseline, but last-layer-only is slightly worse. Lower loss
    is better.}
  \label{fig:f4}
\end{figure*}

We deliberately ablate on \emph{text modeling} --- the stable signal identified in Part 1 ---
to isolate configuration effects from the seed variance that contaminates synthetic-task
comparisons. On text, \Canon-ABCD reaches a final loss of $0.781$ versus $0.793$ for the \Llama
baseline, and \Canon-AC (positions A and C only) reaches $0.782$: positions B and D contribute
negligibly. The remaining ablations (Figure~\ref{fig:f4}) are consistent and interpretable.
\textbf{Kernel size:} $k=1$ reduces to a pointwise rescaling and is worse than no \Canon
($0.797$ vs.\ $0.793$), while $k\ge2$ improves monotonically ($k{=}6$: $0.778$); cross-token
mixing is essential. \textbf{Initialization:} all schemes, including zero init, beat the
baseline, with the $1/k$ constant scheme best. \textbf{Layer depth:} all-layers is strongest,
first-layer-only still beats the baseline, and last-layer-only is slightly worse --- inserting
\Canon only at the final layer is counterproductive. The benefit is also normalization-agnostic,
holding under Pre-LN, Post-LN, and Peri-LN.

We emphasize that the text effect is small in absolute terms ($\approx0.01$), so claims that
distinguish \emph{adjacent} configurations require multi-seed error bars; establishing those
intervals is ongoing work. Configuration tells us how to deploy \Canon but not where its benefit
comes from. For that we decompose by scale and task.

\subsection{Where the Gains Live}
\label{sec:p2:where}

\begin{figure*}[t]
  \centering
  \begin{subfigure}[b]{0.46\linewidth}
    \includegraphics[width=\linewidth]{loss_scaling.png}
    \caption{Aggregate loss vs.\ scale.}
    \label{fig:f5:scale}
  \end{subfigure}\hfill
  \begin{subfigure}[b]{0.46\linewidth}
    \includegraphics[width=\linewidth]{per_task_breakdown_30m.png}
    \caption{Per-task \Canon$-$\Llama delta at 30M.}
    \label{fig:f5:task}
  \end{subfigure}
  \caption{\textbf{The gains are scale-local and task-local.}
    (a) The \Canon$-$\Llama gap shrinks monotonically with scale ($+0.086$ at 3M to $+0.023$ at
    100M) and is within seed noise by 30M (Figure~\ref{fig:f3}). (b) At 30M the improvement is
    concentrated almost entirely in \Depo (multi-hop retrieval, $+0.57$ to $+0.71$); tasks near
    ceiling (\Bfs, \ShortPath) gain little or nothing. A large \Depo accuracy gain coexists with
    a null aggregate because the aggregate is dominated by near-ceiling tasks.}
  \label{fig:f5}
\end{figure*}

\paragraph{Scale.} Across 3M, 10M, 30M, and 100M parameters (five seeds each), the \Canon$-$\Llama
loss gap shrinks monotonically: $+0.086$ at 3M, $+0.045$ at 10M, $+0.031$ at 30M, $+0.023$ at
100M (Figure~\ref{fig:f5:scale}). Combined with the seed-flip audit of
Section~\ref{sec:p1:flip}, which shows the 30M point is within noise, the honest claim is that
\Canon is significantly better \emph{at small scale} and that the benefit dissolves into seed
noise by 30M --- not that \Canon is uniformly better. Applying the seed-robustness audit at each
scale to confirm where the gap is statistically separable from zero is ongoing.

\paragraph{Per task.} At 30M the per-task decomposition (Figure~\ref{fig:f5:task}) shows the
improvement is concentrated almost entirely in \Depo: $+0.710$ on the adjacency encoding and
$+0.574$ on the edge encoding, against deltas between $-0.007$ and $+0.062$ on the other six
tasks. This itself argues for decomposition over aggregation: a large \Depo gain barely moves
the aggregate, because the aggregate is dominated by tasks at the accuracy ceiling
(\Bfs, \ShortPath at $0.91$--$0.97$). The concentration is, however, confounded: \Llama is at the
\emph{floor} on \Depo and at the \emph{ceiling} on \Bfs/\ShortPath, so ``\Canon helps \Depo
specifically'' cannot yet be separated from ``\Canon helps wherever there is headroom.''
Disentangling this requires re-running \Bfs/\ShortPath at higher difficulty so the baseline sits
mid-range (a de-ceiling experiment); this is the highest-value outstanding run. Wherever the
robust effect lives, we finally ask what mechanism produces it.

\subsection{Why: Control of Deep-Layer Gradient Dominance}
\label{sec:p2:mechanism}

\begin{figure*}[t]
  \centering
  \begin{subfigure}[b]{0.32\linewidth}
    \includegraphics[width=\linewidth]{gradient_comparison_synth_vs_text.png}
    \caption{Gradient ratio over training.}
    \label{fig:f6:grad}
  \end{subfigure}\hfill
  \begin{subfigure}[b]{0.32\linewidth}
    \includegraphics[width=\linewidth]{canon_weights_synthetic.png}
    \caption{Learned \Canon weights by shift.}
    \label{fig:f6:weights}
  \end{subfigure}\hfill
  \begin{subfigure}[b]{0.32\linewidth}
    \includegraphics[width=\linewidth]{kurtosis_comparison_synth_vs_text.png}
    \caption{Activation kurtosis by layer.}
    \label{fig:f6:kurt}
  \end{subfigure}
  \caption{\textbf{What \Canon does to training dynamics.} (a) \Llama develops strong deep-layer
    gradient dominance over training (deep/shallow ratio rising to $2.8\times$ on text,
    $1.72\times$ on synthetic), whereas \Canon holds the ratio well below one ($0.70\times$ text,
    $0.44\times$ synthetic); both start near $0.5\times$. (b) \Canon weights favor farther shifts
    from the first checkpoint, a structural bias toward long-range integration. (c) \Canon
    reshapes activation statistics in a domain-dependent way --- suppressing kurtosis on text,
    elevating it on synthetic --- which we report descriptively, not as the performance
    mechanism.}
  \label{fig:f6}
\end{figure*}

We caveat this section up front: the dynamics runs are at 30M, where the aggregate effect is null
(Section~\ref{sec:p1:flip}). We therefore characterize \emph{what \Canon does to training
dynamics} and ask whether it plausibly translates into performance, rather than assuming \Canon
``wins'' at 30M.

The candidate mechanism is control of the deep/shallow gradient ratio. \Llama develops pronounced
deep-layer gradient dominance over training --- the ratio rises from $\approx0.5\times$ at the
start to $2.8\times$ on text and $1.72\times$ on synthetic --- while \Canon holds the ratio well
below one ($0.70\times$ text, $0.44\times$ synthetic) throughout (Figure~\ref{fig:f6:grad}). For
precision we note that on synthetic tasks \Canon's $0.44\times$ is in fact \emph{further} from
uniform than \Llama's $1.72\times$; the universal pattern is a downward shift that suppresses
deep-layer dominance, not a move toward uniformity. The learned weights $\{w_i\}$ in
\eqref{eq:canon} favor farther shifts ($w_2 > w_0$) from the very first checkpoint
(Figure~\ref{fig:f6:weights}), so \Canon is structurally biased toward integrating non-local
context from initialization.

Whether this gradient-flow control is causal for performance remains to be established. The
load-bearing test is a depth ablation at fixed parameter count: if the \Canon$-$\Llama performance
gap grows with depth --- where gradient pathologies worsen --- the mechanism matters for
performance; if the gap is flat, the dynamics are a correlate rather than a cause. This test is
ongoing. Finally, \Canon reshapes activation kurtosis in a domain-dependent way --- suppressing it
on text (\Llama L1 $\approx157$ vs.\ \Canon $\approx18$) and elevating it on synthetic tasks
(\Canon L4 $\approx14.7$ vs.\ \Llama $\approx2.0$; Figure~\ref{fig:f6:kurt}). Because both directions
coincide with improved performance in their respective domains, we report this as a descriptive
reshaping of activation statistics and make no causal performance claim.

% ===========================================================================
\section{Discussion and Implications}
\label{sec:discussion}
% ===========================================================================

\paragraph{Methodological (primary).} Physics-style synthetic reasoning benchmarks demand
\emph{multi-seed, multi-task, multi-scale} evaluation, and verdict stability should be reported as
a function of seed count. A single aggregate number at a single scale is untrustworthy: it
compounds seed noise (Section~\ref{sec:p1:variance}) with ceiling (\Bfs/\ShortPath) and floor
(\Depo-in-mix) effects (Section~\ref{sec:p2:where}). We exhibit a verdict that is null and flips
under exactly the single-number protocol the field currently uses --- the strongest available
argument for the protocol we advocate.

\paragraph{On \Canon.} The benefit is real but conditional: robust at small scale and on
multi-hop retrieval (\Depo), null in the 30M aggregate, and associated with control of deep-layer
gradient dominance (causal status pending the depth test). Practitioners should adopt \Canon with
this scope in mind rather than as a universal improvement, and should prefer the AC configuration
with $k\ge2$ across all layers.

\paragraph{Limitations.} Our configuration ablations are run on text, where the effect is small,
and multi-seed intervals on adjacent configurations are still being collected. The largest scale
tested is 100M. The mechanism analysis rests on a depth test rather than a direct gradient-flow
intervention, and the dynamics runs sit at a scale with no aggregate effect. The de-ceiling and
depth experiments needed to fully close the per-task and mechanism arguments are in progress.

\paragraph{Future work.} Direct intervention on the gradient ratio would upgrade the mechanism
from association to cause; comparing fixed- versus learned-weight \Canon would isolate the
structural inductive bias from the learned one; and extending the scaling study beyond 100M would
test whether the small-scale advantage truly vanishes or merely shrinks below current measurement
resolution.

% ===========================================================================
\begin{thebibliography}{9}
\bibitem{allen_zhu_physics4}
  Z.\ Allen-Zhu and Y.\ Li.
  \emph{Physics of Language Models: Part 4.1 --- Architecture Design and the Magic of Canon
  Layers}.
  arXiv preprint, 2025.

\bibitem{power2022grokking}
  A.\ Power, Y.\ Burda, H.\ Edwards, I.\ Babuschkin, and V.\ Misra.
  \emph{Grokking: Generalization Beyond Overfitting on Small Algorithmic Datasets}.
  arXiv:2201.02177, 2022.
\end{thebibliography}

\end{document}
